\section{Unified Utility: Baselines, Method, and Variants}
\label{sec:utility-results}

\subsection{Legacy vs.\ unified utility: not the same scale}

Two utility functions are implemented, selected by \texttt{--utility
legacy|unified}. \textbf{Legacy} is a single weighted quality-times-freshness
sum with no tardiness, cost, or fairness term, using a pixel-fidelity
quality proxy ($1-\text{LPIPS}$). \textbf{Unified} adds a tardiness penalty,
a resource-cost penalty, and a maximin fairness floor, and replaces the
proxy with the measured downstream mIoU term $s_q$ from
\S\ref{sec:downstream} (full formulation in \S\ref{sec:formulation}). These
two modes are \emph{not} on the same scale and must not be compared directly
in a single table; Table~\ref{tab:legacy-vs-unified} reports each mode's
own best-fixed-depth-greedy vs.\ MPC margin, from the \texttt{gbs-limited}
scenario, to show how switching the quality term changes the conclusion.

\begin{table}[h]
\centering
\caption{MPC vs.\ best fixed-depth greedy, gbs-limited scenario, seed 42.}
\label{tab:legacy-vs-unified}
\begin{tabular}{lrrr}
\toprule
Utility mode & \texttt{mpc} & best fixed greedy & MPC margin \\
\midrule
legacy ($1-\text{LPIPS}$ quality) & 64.95 & 64.15 (fixed-8) & $+1.2\%$ \\
unified (real mIoU quality)       & 18.13 & 13.71 (fixed-8) & $+32.2\%$ \\
\bottomrule
\end{tabular}
\end{table}

Under the pixel-fidelity proxy, MPC essentially ties the best fixed-depth
greedy baseline. Once quality reflects the real downstream task, MPC's
percentage margin is $\approx 26\times$ larger. This is the central result
motivating the unified utility: the legacy proxy under-states MPC's value
because it is nearly flat across compression depths (\S\ref{sec:downstream}).

\subsection{Compared algorithms}

\begin{itemize}
\item \textbf{greedy-fixed-\{2,4,8,16\}}: always encode at a fixed depth,
  transmit whatever fits in the current slot's remaining capacity.
\item \textbf{greedy-adaptive}: greedy, but adapts encode depth to queue
  length.
\item \textbf{mpc}: flat rolling-horizon MILP (\S\ref{sec:formulation}),
  static shortest-path routing.
\item \textbf{mpc-congestion}: flat MPC with predictive, congestion-aware
  routing (M/M/1-style edge cost, iterated to a fixed point).
\item \textbf{mpc-hier}\todoconfirm{scope of this variant to keep/drop is
  pending Xuanhao's confirmation, see note below}: two-level hierarchy —
  a slow admission/budget LP plus a fast per-task depth MILP, with explicit
  task drop/abandon logic.
\item \textbf{mpc-hier-route}: \texttt{mpc-hier} plus a routing MPC that
  freezes a concrete path per 10-minute macro-epoch.
\item \textbf{mpc-2level}: a routing MPC (multicommodity-flow LP over
  candidate paths) and the depth MPC iterated to a fixed point as
  \emph{peers}, not hierarchically.
\end{itemize}

\todoconfirm{Xuanhao's email asked to drop ``the previous mpc-hier results
where the two MPCs are both used for depth selection.'' \texttt{FORMULATION.md}
describes \texttt{mpc-hier}'s two levels as: upper = LP over a
\emph{continuous relaxed depth} $x_{k,q}\in[0,1]$ (for budgeting), lower =
MILP over \emph{integral depth} $x_{k,q}$ -- i.e.\ both levels of
\texttt{mpc-hier} specifically operate on depth, with static routing at
both. This is different from \texttt{mpc-hier-route}/\texttt{mpc-2level},
which are an explicit routing-MPC + depth-MPC split. This textual match
makes \texttt{mpc-hier} the likely candidate for what Xuanhao means, but all
schedulers are reported below pending his direct confirmation --
see \S\ref{sec:hier} for the full two-level description.}

\subsection{Ten schedulers, one run}

Table~\ref{tab:scheduler-comparison} reports all ten schedulers on the
\texttt{fabric-limited} scenario (where routing choices are measurable),
unified utility, real mIoU quality table, single run/seed (42). Every row is
from the same run, so utility, delivery, on-time rate, and bound-gap are
directly comparable across schedulers.

\begin{table}[h]
\centering
\caption{Scheduler comparison, fabric-limited scenario, seed 42. Bound
(offline HiGHS upper bound, \S\ref{sec:formulation}) $= 18.99$.}
\label{tab:scheduler-comparison}
\begin{tabular}{lrrrr}
\toprule
Scheduler & Utility & Delivery\% & On-time\% & Gap to bound \\
\midrule
greedy-fixed-2   & 6.43  & 100\% & 88.7\% & 66.1\% \\
greedy-fixed-4   & 11.08 & 100\% & 88.7\% & 41.7\% \\
greedy-fixed-8   & 13.71 & 100\% & 88.6\% & 27.8\% \\
greedy-fixed-16  & 8.80  & 100\% & 81.5\% & 53.7\% \\
greedy-adaptive  & 8.80  & 100\% & 81.5\% & 53.7\% \\
mpc              & 18.76 & 87.9\% & 86.3\% & 1.2\% \\
mpc-congestion   & 18.86 & 89.3\% & 87.6\% & 0.7\% \\
mpc-hier         & 14.76 & 78.7\% & 75.0\% (39 dropped) & 22.3\% \\
mpc-hier-route   & 14.78 & 81.9\% & 78.6\% (40 dropped) & 22.2\% \\
\textbf{mpc-2level} & \textbf{18.90} & \textbf{90.3\%} & \textbf{88.6\%} & \textbf{0.5\%} \\
\bottomrule
\end{tabular}
\end{table}

\emph{Note:} \texttt{greedy-adaptive}'s row is numerically identical to
\texttt{greedy-fixed-16} in this run — in this scenario, the adaptive
policy's depth choice converges to always selecting depth 16, so this is a
genuine coincidence in the underlying decisions (verified in
\S\ref{sec:decisions}), not a duplicated row.

\subsection{Individual utility components}

Table~\ref{tab:utility-components} decomposes total utility into its four
additive terms, from the same run as Table~\ref{tab:scheduler-comparison}
(\texttt{hypatia\_sim/oec\_scenario\_couplings/summary.txt}).
$U = \text{quality} - \text{tardy} - \text{cost} + \text{fair}$; $u_{\min}$
is the optimized maximin fairness floor and Jain's index is reported as a
diagnostic only (not MILP-representable, so not itself optimized; see
\S\ref{sec:formulation}).

\begin{table}[h]
\centering
\caption{Utility component decomposition, fabric-limited scenario, seed 42.}
\label{tab:utility-components}
\begin{tabular}{lrrrrrrrr}
\toprule
Scheduler & Total & Quality & $-$Tardy & $-$Cost & $+$Fair & $u_{\min}$ & Jain & Enc.\ (kJ) \\
\midrule
mpc              & 18.757 & 16.344 & $-$0.029 & $-$2.121 & 4.564  & 0.7131  & 0.9978 & 4170.2 \\
greedy-adaptive  & 8.797  & 16.170 & $-$1.502 & $-$2.350 & $-$3.521 & $-$0.5502 & 0.9118 & 4742.3 \\
greedy-fixed-2   & 6.431  & 8.414  & $-$0.813 & $-$1.322 & 0.152  & 0.0237  & 0.9331 & 4742.3 \\
greedy-fixed-4   & 11.076 & 12.089 & $-$0.813 & $-$1.469 & 1.269  & 0.1983  & 0.9670 & 4742.3 \\
greedy-fixed-8   & 13.711 & 14.394 & $-$0.814 & $-$1.763 & 1.894  & 0.2959  & 0.9748 & 4742.3 \\
greedy-fixed-16  & 8.797  & 16.170 & $-$1.502 & $-$2.350 & $-$3.521 & $-$0.5502 & 0.9118 & 4742.3 \\
mpc-congestion   & 18.864 & 16.447 & $-$0.028 & $-$2.151 & 4.597  & 0.7182  & 0.9985 & 4235.1 \\
mpc-hier         & 14.761 & 13.743 & $-$0.085 & $-$1.536 & 2.639  & 0.4124  & 0.9846 & 3733.3 \\
\textbf{mpc-2level} & \textbf{18.898} & \textbf{16.503} & $-$0.030 & $-$2.173 & 4.597 & \textbf{0.7183} & 0.9987 & 4281.9 \\
mpc-hier-route   & 14.779 & 13.811 & $-$0.055 & $-$1.561 & 2.584  & 0.4038  & 0.9848 & 3885.9 \\
\bottomrule
\end{tabular}
\end{table}

\textbf{What drives the gap: admission control, not routing.}
\texttt{mpc-2level} wins outright, but \texttt{mpc-congestion} (predictive
routing only, no full flow-optimization) sits within 0.04 utility of it. A
five-point ISL-rate sweep (\S\ref{sec:sweeps}) found the dedicated
routing-LP's own marginal contribution tops out at $+0.84\%$ utility and is
exactly zero above 1.5\,Mbps ISL. The real, structural $\approx 5$-point gap
is between the no-admission-control family (\texttt{mpc},
\texttt{mpc-congestion}, \texttt{mpc-2level}, $\approx 18.8$--$18.9$) and the
admission-control family (\texttt{mpc-hier}, \texttt{mpc-hier-route},
$\approx 14.8$) — whether to drop a task early matters far more than how
cleverly the kept tasks are routed.

\textbf{Route-freezing does not survive a LEO fabric.} A route frozen for a
10-minute macro-epoch (\texttt{mpc-hier-route}) is usable at its
actually-executed slot only 7.7\% of the time, because ground-station
contact windows (227--265\,s, Table~\ref{tab:constellation}) are shorter
than any freeze period worth having. This is a structural property of the
constellation geometry, not a tuning failure.

\subsection{Solve-time and energy efficiency}

\texttt{mpc-hier}/\texttt{mpc-hier-route} solve $3$--$4\times$ faster per
call (mean 43.5--43.6\,ms) than the flat/two-level MPCs (123.9--168.5\,ms
mean, up to 3{,}088\,ms max for \texttt{mpc-2level}), at the cost of the
utility gap above. \texttt{mpc-2level} solves 365 times over the run
(vs.\ 113--121 for the other MPC variants), for 45.2\,s total solve time
vs.\ 20.4\,s for flat \texttt{mpc} — recall from \S\ref{sec:setup} that these
are individual-solve statistics summed, not simulated wall-clock time.
